\input{../tex-inputs/latex-korrekturansicht-vorspann}

\section[Ludwig Fulda u. a. an Arthur Schnitzler, 18. 9. 1902]{L04348 Ludwig Fulda u. a. an Arthur Schnitzler, 18. 9. 1902}
\nopagebreak\mylabel{L04348v}
\rehead{ }\normalsize\beginnumbering\briefempfaengerindex{Schnitzler, Arthur@\textsc{Schnitzler, Arthur}!zzzLeipziger-Stettenheim, Julie@\emph{von Julie Leipziger-Stettenheim}!1902-09-181@{{[}18. 9. 1902{]}}|(be}\briefempfaengerindex{Schnitzler, Arthur@\textsc{Schnitzler, Arthur}!zzzBerény, Henri@\emph{von Henri Berény}!1902-09-181@{{[}18. 9. 1902{]}}|(be}\briefempfaengerindex{Schnitzler, Arthur@\textsc{Schnitzler, Arthur}!zzzWiehe, Charlotte@\emph{von Charlotte Wiehe}!1902-09-181@{{[}18. 9. 1902{]}}|(be}\briefempfaengerindex{Schnitzler, Arthur@\textsc{Schnitzler, Arthur}!zzzd’Albert, Ida@\emph{von Ida d’Albert}!1902-09-181@{{[}18. 9. 1902{]}}|(be}\briefempfaengerindex{Schnitzler, Arthur@\textsc{Schnitzler, Arthur}!zzzGoldmann, Paul@\emph{von Paul Goldmann}!1902-09-181@{{[}18. 9. 1902{]}}|(be}\briefempfaengerindex{Schnitzler, Arthur@\textsc{Schnitzler, Arthur}!zzzFulda, Ludwig@\emph{von Ludwig Fulda}!1902-09-181@{{[}18. 9. 1902{]}}|(be}
\toendnotes[C]{\smallbreak\pagebreak[2]}
\correspDesc{Versand  durch Ludwig Fulda, Paul Goldmann, Ida d’Albert, Charlotte Wiehe, Henri Berény, Julie Leipziger-Stettenheim am [18. 9. 1902] in Berlin
\newline{}Übermittlung  am 18. 9. 1902 in Berlin
\newline{}Zustellung  am 19. 9. 1902 in Wien
\newline{}Erhalt  durch Arthur Schnitzler am [19. 9. 1902] in Wien}\toendnotes[C]{\smallbreak}
\Standort{CUL, Schnitzler, B 331.}
\physDesc{Postkarte, 370 Zeichen
\newline{}Handschrift Ida d’Albert: blaue Tinte, lateinische Kurrent
\newline{}Handschrift Charlotte Wiehe: schwarze Tinte, lateinische Kurrent
\newline{}Handschrift Paul Goldmann: schwarze Tinte, deutsche Kurrent
\newline{}Handschrift Henri Berény: schwarze Tinte, lateinische Kurrent
\newline{}Handschrift Julie Leipziger-Stettenheim: schwarze Tinte, lateinische Kurrent
\newline{}Handschrift Ludwig Fulda: schwarze Tinte, deutsche Kurrent
\newline{}Versand: 1) Stempel: »\nobreak{}\oindex{Charlottenburg@\textbf{Charlottenburg}, \emph{Ehemaliger Ort}|pwk}\textcolor{gray}{Ch}arlottenburg, 18. 9. {[}1902{]}, 8–9\textcolor{gray}{N}\nobreak{}«.   2) Stempel: »\nobreak{}\oindex{IX., Alsergrund@\textbf{IX., Alsergrund}, \emph{Verwaltungsgebiet}|pwk}9/3 {[}Wien{]}, 19. 9. 02, 7.N, Best{[}ellt{]}\nobreak{}«. }\toendnotes[C]{\smallbreak}\pstart{}{\pb}{[}hs. Fulda:{]} Herrn\pend{}\pstart{}Dr. \textsc{Arthur Schnitzler}\pend{}\pstart{}\textsc{\textcolor{pink}{Wien}\oindex{Wien@\textbf{Wien}, \emph{Verwaltungsgebiet}|pw}{}\ledrightnote{\textcolor{pink}{Wien}}}\pend{}\pstart{}\textsc{\textcolor{pink}{IX. Frangasse 1}\oindex{Wien@\textbf{Wien}!IX., Alsergrund@\textbf{IX., Alsergrund}!Frankgasse 1@\textbf{Frankgasse 1}, \emph{Wohngebäude}|pw}{}\ledrightnote{\textcolor{pink}{Frankgasse 1}}.}\pend{}{\bigskip}\vspace{1em}
\pstart
           \noindent{}{\pb}{[}hs. d’Albert:{]} In Gesellschaft Ihrer entzückenden \label{K_L04348-1v}\edtext{französischen \textcolor{green}{Annie}\pwindex{Schnitzler, Arthur 15. 5. 1862 Wien – 21. 10. 1931 ebd.@\textsc{Schnitzler, Arthur} (15. 5. 1862 Wien – 21. 10. 1931 ebd.), \emph{Schriftsteller, Mediziner}!Souper d’Adieu@\strich\emph{Souper d’Adieu}|pwv}{}\ledrightnote{{$\rightarrow$}\emph{\textcolor{green}{Souper d’Adieu}}}}{\lemma{\textnormal{\emph{französischen Annie}}}\Cendnote{\textnormal{\textcolor{blue}{Charlotte Wiehe}\pwindex{Wiehe, Charlotte 28.\,8.\,1865 Kopenhagen – 4.\,9.\,1947 Vedbæk@\textsc{Wiehe, Charlotte} (28.\,8.\,1865 Kopenhagen – 4.\,9.\,1947 Vedbæk), \emph{Schauspielerin}|pwk} tourte
                  gerade und hatte \emph{\textcolor{green}{Souper d’Adieu}\pwindex{Schnitzler, Arthur 15. 5. 1862 Wien – 21. 10. 1931 ebd.@\textsc{Schnitzler, Arthur} (15. 5. 1862 Wien – 21. 10. 1931 ebd.), \emph{Schriftsteller, Mediziner}!Souper d’Adieu@\strich\emph{Souper d’Adieu}|pwk}} (\emph{\textcolor{green}{Abschiedssouper}\pwindex{Schnitzler, Arthur 15. 5. 1862 Wien – 21. 10. 1931 ebd.@\textsc{Schnitzler, Arthur} (15. 5. 1862 Wien – 21. 10. 1931 ebd.), \emph{Schriftsteller, Mediziner}!Abschiedssouper@\strich\emph{Abschiedssouper}|pwk}}) im Repertoire. \textcolor{blue}{Schnitzler} sah eine Aufführung am 29. 9. 1902 im \textcolor{pink}{Raimund-Theater}\oindex{Wien@\textbf{Wien}!VI., Mariahilf@\textbf{VI., Mariahilf}!Raimund-Theater@\textbf{Raimund-Theater}, \emph{Theater}|pwk}. }}}\label{K_L04348-1}\pend
           \pstart grüssen schönstens Ihre \spacefill\mbox{Idus Fulda}\pend{}\selectlanguage{ngerman}\vspace{1em}
\pstart
           \noindent{}{[}hs. Wiehe:{]} Mes meilleurs compliments!\pend
           \pstart \spacefill\mbox{Charlotte Wiehe-Berény}\pend{}\selectlanguage{ngerman}\vspace{1em}
\pstart
           \noindent{}{[}hs. Leipziger-Stettenheim:{]} Wann werde ich Sie endlich kennen lernen? Zeit wäre
               es schon!\pend
           \pstart \spacefill\mbox{Julie Leipziger.}\pend{}\selectlanguage{ngerman}\vspace{1em}
\pstart
           \noindent{}{[}hs. Berény:{]} Ich hoffe Sie \label{K_L04348-2v}\edtext{bald in \textcolor{pink}{Wien}\oindex{Wien@\textbf{Wien}, \emph{Verwaltungsgebiet}|pw}{}\ledrightnote{\textcolor{pink}{Wien}}}{\lemma{\textnormal{\emph{bald in Wien}}}\Cendnote{\textnormal{Vgl. A. S.: \emph{Tagebuch}, 28. 9. 1902.}}}\label{K_L04348-2} begrüssen zu können.\pend
           \pstart Henri Bereny\pend{}\selectlanguage{ngerman}\vspace{1em}
\pstart
           \noindent{}{[}hs. Goldmann:{]} Herzlichen Gruß!\pend
           \pstart \spacefill\mbox{Paul Goldmann.}\pend{}\selectlanguage{ngerman}\vspace{1em}
\pstart
           \noindent{}{[}hs. Fulda:{]} Und auch von mir!\pend
           \pstart \spacefill\mbox{Ludwig Fulda}\pend{}\selectlanguage{ngerman}\endnumbering\briefempfaengerindex{Schnitzler, Arthur@\textsc{Schnitzler, Arthur}!zzzLeipziger-Stettenheim, Julie@\emph{von Julie Leipziger-Stettenheim}!1902-09-181@{{[}18. 9. 1902{]}}|)be}\briefempfaengerindex{Schnitzler, Arthur@\textsc{Schnitzler, Arthur}!zzzBerény, Henri@\emph{von Henri Berény}!1902-09-181@{{[}18. 9. 1902{]}}|)be}\briefempfaengerindex{Schnitzler, Arthur@\textsc{Schnitzler, Arthur}!zzzWiehe, Charlotte@\emph{von Charlotte Wiehe}!1902-09-181@{{[}18. 9. 1902{]}}|)be}\briefempfaengerindex{Schnitzler, Arthur@\textsc{Schnitzler, Arthur}!zzzd’Albert, Ida@\emph{von Ida d’Albert}!1902-09-181@{{[}18. 9. 1902{]}}|)be}\briefempfaengerindex{Schnitzler, Arthur@\textsc{Schnitzler, Arthur}!zzzGoldmann, Paul@\emph{von Paul Goldmann}!1902-09-181@{{[}18. 9. 1902{]}}|)be}\briefempfaengerindex{Schnitzler, Arthur@\textsc{Schnitzler, Arthur}!zzzFulda, Ludwig@\emph{von Ludwig Fulda}!1902-09-181@{{[}18. 9. 1902{]}}|)be}\mylabel{L04348h}
\begin{anhang}
\end{anhang}\newcommand{\dateiname}{L04348}\newcommand{\titel}{Ludwig Fulda u. a. an Arthur Schnitzler, 18. 9. 1902}\newcommand{\editorInnen}{Herausgegeben von Jahnke, SelmaMüller, Martin Anton}\input{../tex-inputs/latex-korrekturansicht-abspann}
